
Economic models are commonly used to simplify and represent complex real world relationships.
Since all models are approximations, additional components are often introduced to better capture
underlying patterns. Especially in micro economices where the assumption of independent and identically
distributed (i.i.d) can not be reasonably assumed due to its time factor or close proximity making
them influence each other.

Spatial models are not a novel concept. This research seeks to fill the
gap in spatial research and identifying an important difficulty in modeling spatial relationships.
Through simulation, this research seeks to understand the limitations of various spatial models
and identify contexts in which specific models are preferable. More specifically in this research
we seek to find reasonable strategies to best specify spatial models, where are it can be reasonably
be presumed that there is some spatial effect but there is no clear answer on how spatial agents relate to each other.
